\documentclass[12pt,a4paper]{article}

% ── Packages ────────────────────────────────────────────────────────────────
\usepackage[T1]{fontenc}
\usepackage[utf8]{inputenc}
\usepackage{lmodern}
\usepackage{microtype}
\usepackage{amsmath,amssymb,amsthm}
\usepackage{mathtools}
\usepackage{enumitem}
\usepackage{geometry}
\usepackage{hyperref}
\usepackage{booktabs}
\usepackage{array}
\usepackage{xcolor}
\usepackage{listings}
\usepackage{fancyhdr}

\geometry{margin=2.5cm}
\setlength{\headheight}{13.60001pt}
\addtolength{\topmargin}{-1.60001pt}

% ── Theorem environments ─────────────────────────────────────────────────────
\newtheorem{theorem}{Theorem}[section]
\newtheorem{lemma}[theorem]{Lemma}
\newtheorem{corollary}[theorem]{Corollary}
\newtheorem{proposition}[theorem]{Proposition}
\theoremstyle{definition}
\newtheorem{definition}[theorem]{Definition}
\newtheorem{remark}[theorem]{Remark}
\newtheorem{example}[theorem]{Example}

% ── Code listings ────────────────────────────────────────────────────────────
\definecolor{codegray}{rgb}{0.96,0.96,0.96}
\definecolor{codeblue}{rgb}{0.1,0.2,0.6}
\definecolor{codegreencomment}{rgb}{0.13,0.55,0.13}

\lstset{
  backgroundcolor=\color{codegray},
  basicstyle=\ttfamily\small,
  keywordstyle=\color{codeblue}\bfseries,
  commentstyle=\color{codegreencomment}\itshape,
  breaklines=true,
  frame=single,
  framerule=0pt,
  rulecolor=\color{gray!30},
  xleftmargin=6pt, xrightmargin=6pt,
}

\lstdefinelanguage{Lean4}{
  keywords={theorem,lemma,def,axiom,by,have,obtain,exact,intro,rw,push_cast,
            ring,decide,linarith,nlinarith,rcases,calc,simp,fin_cases,revert,
            import,where,fun,if,then,else,let,in,match,with,show,apply,refine,
            constructor,left,right,norm_num,norm_cast,exact_mod_cast,
            unfold,push_neg,omega,Set,Finset,open,use,private,rintro,neg_inj},
  sensitive=true,
  morecomment=[l]{--},
  morecomment=[s]{/-}{-/},
  morestring=[b]",
  keywordstyle=\color{codeblue}\bfseries,
  commentstyle=\color{codegreencomment}\itshape,
}

% ── Macros ───────────────────────────────────────────────────────────────────
\newcommand{\Z}{\mathbb{Z}}
\newcommand{\Q}{\mathbb{Q}}
\newcommand{\N}{\mathbb{N}}

% ── Header / Footer ──────────────────────────────────────────────────────────
\pagestyle{fancy}
\fancyhf{}
\rhead{\small\thepage}
\lhead{\small Infinitely Many Solutions to $x^5 + y^4 + 3z^2 = 0$}
\renewcommand{\headrulewidth}{0.4pt}

% ── Title ────────────────────────────────────────────────────────────────────
\title{%
  \Large\bfseries Infinitely Many Integer Solutions to\\[6pt]
  \Huge $x^5 + y^4 + 3z^2 = 0$\\[10pt]
  \large Complete Analysis with Lean~4 Formalisation
}
\author{%
  \normalsize Verified in Lean~4 / Mathlib~4 (v4.21.0)\\[2pt]
  \normalsize Repository: \href{https://github.com/JAgbanwa/miscellaneousmathproblems}%
    {\texttt{JAgbanwa/miscellaneousmathproblems}}
}
\date{\normalsize April 2026}

% ════════════════════════════════════════════════════════════════════════════
\begin{document}

\maketitle
\tableofcontents
\newpage

% ════════════════════════════════════════════════════════════════════════════
\section{Statement and Main Result}\label{sec:intro}

\subsection{The equation}

We study the Diophantine equation
\begin{equation}\label{eq:main}
  x^5 + y^4 + 3z^2 = 0, \qquad (x,y,z)\in\Z^3. \tag{$\star$}
\end{equation}

\subsection{Main theorem}

\begin{theorem}[Main result]\label{thm:main}
  The equation $x^5 + y^4 + 3z^2 = 0$ has \textbf{infinitely many} integer solutions.
  In particular, the parametric family
  \begin{equation}\label{eq:family1}
    (x, y, z) = \bigl(-n^4,\; \varepsilon\, n^5,\; 0\bigr),
    \qquad n \in \Z,\; \varepsilon \in \{+1,-1\},
  \end{equation}
  provides infinitely many distinct integer solutions.  A second independent family is
  \begin{equation}\label{eq:family2}
    (x, y, z) = \bigl(-3k^2,\; 0,\; \varepsilon\, 9k^5\bigr),
    \qquad k \in \Z,\; \varepsilon \in \{+1,-1\},
  \end{equation}
  and a third is
  \begin{equation}\label{eq:family3}
    (x, y, z) = \bigl(-4t^4,\; \varepsilon\, 4t^5,\; \delta\, 16t^{10}\bigr),
    \qquad t \in \Z,\; \varepsilon,\delta \in \{+1,-1\}.
  \end{equation}
\end{theorem}

The proofs are elementary: each family is verified by a single ring identity,
and injectivity is established by strict monotonicity of $n \mapsto n^4$ on $\N$.
See Section~\ref{sec:derivation} for derivations,
Section~\ref{sec:proof} for the complete proof,
and Section~\ref{sec:lean} for the Lean~4 formalisation.

% ════════════════════════════════════════════════════════════════════════════
\section{Sign Constraint}\label{sec:sign}

Since $y^4 \geq 0$ and $3z^2 \geq 0$ for all $(y,z)\in\Z^2$, equation~\eqref{eq:main}
requires $x^5 \leq 0$, hence $x \leq 0$.  Writing $x = -a$ with $a \geq 0$
transforms~\eqref{eq:main} into the equivalent form
\begin{equation}\label{eq:reduced}
  y^4 + 3z^2 = a^5, \quad a \geq 0;\qquad (a,y,z)\in\N_0 \times \Z^2. \tag{$\star\star$}
\end{equation}
It therefore suffices to produce infinitely many integer solutions to~\eqref{eq:reduced}.

% ════════════════════════════════════════════════════════════════════════════
\section{Derivation of the Three Families}\label{sec:derivation}

\subsection{Family 1: setting \texorpdfstring{$z = 0$}{z = 0}}

Setting $z = 0$ in~\eqref{eq:main} gives $x^5 + y^4 = 0$, i.e.\ $x^5 = -y^4$.
Since $y^4 \geq 0$ we need $x \leq 0$.  The only integer solutions are
\[
  x = -n^4, \quad y = \pm n^5, \qquad n \in \Z,
\]
because $(-n^4)^5 = -n^{20}$ and $(n^5)^4 = n^{20}$, so the sum is zero.
These are in fact all solutions with $z = 0$: if $x^5 = -y^4$ and $x = -m$ then
$m^5 = y^4$, forcing $m$ to be a perfect fourth power (say $m = n^4$) and $y = \pm n^5$.

\subsection{Family 2: setting \texorpdfstring{$y = 0$}{y = 0}}

Setting $y = 0$ in~\eqref{eq:main} gives $x^5 = -3z^2$.
Since $3 \mid x^5$ we need $3 \mid x$.  Write $x = -3k^2$ (choosing $x$ as
$-3$ times a perfect square so that $x^5$ acquires the correct algebraic structure):
\[
  (-3k^2)^5 = -3^5 k^{10} = -243k^{10}
  \qquad\text{and}\qquad
  3z^2 = 243k^{10} \;\Longrightarrow\; z^2 = 81k^{10} = (9k^5)^2,
\]
so $z = \pm 9k^5$.  We obtain the family $(-3k^2, 0, 9k^5)$ for all $k \in \Z$.
A brute-force search (see \texttt{brute\_force\_search.py}) confirms this is the
complete family with $y = 0$.

\subsection{Family 3: weighted homogeneity}

Assign \emph{degree weights}
\[
  \mathrm{wt}(x) = 4, \quad \mathrm{wt}(y) = 5, \quad \mathrm{wt}(z) = 10.
\]
Every monomial in~\eqref{eq:main} carries the same weighted degree~20:
\[
  \mathrm{wt}(x^5) = 5 \times 4 = 20, \quad
  \mathrm{wt}(y^4) = 4 \times 5 = 20, \quad
  \mathrm{wt}(3z^2) = 2 \times 10 = 20.
\]

\begin{lemma}[Weighted scaling invariance]\label{lem:scale}
  If $(x_0, y_0, z_0) \in \Z^3$ satisfies~\eqref{eq:main}, then so does
  $(t^4 x_0,\; t^5 y_0,\; t^{10} z_0)$ for every $t \in \Z$.
\end{lemma}

\begin{proof}
  $(t^4 x_0)^5 + (t^5 y_0)^4 + 3(t^{10}z_0)^2
  = t^{20} x_0^5 + t^{20} y_0^4 + 3t^{20} z_0^2
  = t^{20}(x_0^5 + y_0^4 + 3z_0^2) = 0.$ \qed
\end{proof}

A direct calculation confirms that $(-4, 4, 16)$ is a solution:
\[
  (-4)^5 + 4^4 + 3\cdot 16^2 = -1024 + 256 + 768 = 0.
\]
Applying Lemma~\ref{lem:scale} to this seed yields Family~3.

% ════════════════════════════════════════════════════════════════════════════
\section{Proof of the Main Theorem}\label{sec:proof}

\begin{proof}[Proof of Theorem~\ref{thm:main}]
  \textbf{(Verification of Family 1.)}
  \[
    (-n^4)^5 + (n^5)^4 + 3\cdot 0^2 = -n^{20} + n^{20} = 0.\quad\checkmark
  \]

  \textbf{(Verification of Family 2.)}
  \[
    (-3k^2)^5 + 0^4 + 3(9k^5)^2 = -3^5 k^{10} + 3\cdot 81\, k^{10}
    = -243k^{10} + 243k^{10} = 0.\quad\checkmark
  \]

  \textbf{(Verification of Family 3.)}
  \[
    (-4t^4)^5 + (4t^5)^4 + 3(16t^{10})^2 = -1024\,t^{20} + 256\,t^{20} + 768\,t^{20}
    = 0.\quad\checkmark
  \]

  \textbf{(Infinitude.)}  For non-negative integers $n_1 < n_2$, the function
  $n \mapsto n^4$ is strictly increasing on~$\N_0$.  Hence $n_1^4 < n_2^4$, and
  the first components $-n_1^4 \neq -n_2^4$ are distinct.  Therefore
  $(-n_1^4, n_1^5, 0) \neq (-n_2^4, n_2^5, 0)$, so Family~1 yields infinitely
  many distinct integer solutions.
  \qed
\end{proof}

% ════════════════════════════════════════════════════════════════════════════
\section{Explicit Solutions}\label{sec:explicit}

\subsection{Family 1}

\begin{table}[h]
\centering
\caption{Members of Family~1: $(-n^4, n^5, 0)$.}
\label{tab:family1}
\renewcommand{\arraystretch}{1.2}
\begin{tabular}{rrrr}
\toprule
$n$ & $x = -n^4$ & $y = n^5$ & $x^5 + y^4 + 3z^2$ \\
\midrule
$0$  & $0$    & $0$     & $0$ \\
$1$  & $-1$   & $1$     & $0$ \\
$-1$ & $-1$   & $-1$    & $0$ \\
$2$  & $-16$  & $32$    & $0$ \\
$-2$ & $-16$  & $-32$   & $0$ \\
$3$  & $-81$  & $243$   & $0$ \\
\bottomrule
\end{tabular}
\end{table}

\subsection{Family 2}

\begin{table}[h]
\centering
\caption{Members of Family~2: $(-3k^2, 0, 9k^5)$.}
\label{tab:family2}
\renewcommand{\arraystretch}{1.2}
\begin{tabular}{rrrr}
\toprule
$k$ & $x = -3k^2$ & $z = 9k^5$ & $x^5 + y^4 + 3z^2$ \\
\midrule
$0$  & $0$    & $0$      & $0$ \\
$1$  & $-3$   & $9$      & $0$ \\
$-1$ & $-3$   & $-9$     & $0$ \\
$2$  & $-12$  & $288$    & $0$ \\
$3$  & $-27$  & $2187$   & $0$ \\
$4$  & $-48$  & $9216$   & $0$ \\
\bottomrule
\end{tabular}
\end{table}

\subsection{Family 3}

\begin{table}[h]
\centering
\caption{Members of Family~3: $(-4t^4, 4t^5, 16t^{10})$.}
\label{tab:family3}
\renewcommand{\arraystretch}{1.2}
\begin{tabular}{rrrr}
\toprule
$t$ & $x = -4t^4$ & $y = 4t^5$ & $z = 16t^{10}$ \\
\midrule
$0$  & $0$    & $0$    & $0$      \\
$1$  & $-4$   & $4$    & $16$     \\
$-1$ & $-4$   & $-4$   & $16$     \\
$2$  & $-64$  & $128$  & $16384$  \\
\bottomrule
\end{tabular}
\end{table}

% ════════════════════════════════════════════════════════════════════════════
\section{Modular Analysis}\label{sec:modular}

\subsection{Modulo 3}

By Fermat's little theorem, $x^5 \equiv x \pmod{3}$.  Since $y^4 \equiv 0$ or $1
\pmod{3}$ and $3z^2 \equiv 0 \pmod{3}$, the equation reduces to
\[
  x + y^4 \equiv 0 \pmod{3}.
\]
This is satisfied in two classes: (a) $3 \mid x$ and $3 \mid y$; or (b)
$x \equiv 2 \pmod{3}$ and $3 \nmid y$.

\subsection{Modulo 4}

One checks that $x^5 \equiv x \pmod{4}$ for odd $x$, $y^4 \equiv 0$ or $1 \pmod{4}$,
and $3z^2 \equiv 0$ or $3 \pmod{4}$.  All four valid parity combinations
(even--even--even, odd--odd--even, odd--even--odd, even--odd--odd) satisfy the
equation modulo~4, so there is no parity obstruction.

\subsection{No obstruction to infinitely many solutions}

The modular analysis confirms that the two ``obvious'' specialisations $z = 0$
and $y = 0$ each pass all congruence conditions.  Combined with the algebraic
constructions above, this proves the existence of infinitely many solutions
in every residue class compatible with the congruence conditions.

% ════════════════════════════════════════════════════════════════════════════
\section{Lean~4 Formalisation}\label{sec:lean}

The main theorem is fully formalised in Lean~4 / Mathlib~4 (v4.21.0) in the file
\texttt{InfinitelyManySolutions.lean}.  The proof has \textbf{zero} \texttt{sorry}s.

\subsection{Key definitions}

\begin{lstlisting}[language=Lean4]
/-- The Diophantine equation x^5 + y^4 + 3*z^2 = 0. -/
def isolution (x y z : Z) : Prop := x ^ 5 + y ^ 4 + 3 * z ^ 2 = 0

/-- The set of all integer solutions. -/
def solutions : Set (Z x Z x Z) :=
  { p | isolution p.1 p.2.1 p.2.2 }
\end{lstlisting}

\subsection{The three parametric families}

Each family is verified by a single \texttt{ring} call:

\begin{lstlisting}[language=Lean4]
-- Family 1: (-n^4, n^5, 0)
theorem family1 (n : Z) : isolution (-(n ^ 4)) (n ^ 5) 0 := by
  unfold isolution; ring

-- Family 2: (-3*k^2, 0, 9*k^5)
theorem family2 (k : Z) : isolution (-3 * k ^ 2) 0 (9 * k ^ 5) := by
  unfold isolution; ring

-- Family 3: (-4*t^4, 4*t^5, 16*t^10)
theorem family3 (t : Z) : isolution (-4 * t ^ 4) (4 * t ^ 5) (16 * t ^ 10) := by
  unfold isolution; ring
\end{lstlisting}

\subsection{Infinitude}

The map $n \mapsto (-n^4, n^5, 0)$ from $\N$ into the solution set is injective,
proved by the strict monotonicity of $n \mapsto n^4$ on $\N$ and \texttt{neg\_inj}:

\begin{lstlisting}[language=Lean4]
def natMap (n : N) : Z x Z x Z := (-(n : Z) ^ 4, (n : Z) ^ 5, 0)

theorem natMap_injective : Function.Injective natMap := by
  intro n1 n2 h
  simp only [natMap, Prod.mk.injEq] at h
  have h4 : (n1 : Z) ^ 4 = (n2 : Z) ^ 4 := neg_inj.mp h.1
  have h4n : n1 ^ 4 = n2 ^ 4 := by exact_mod_cast h4
  exact pow4_strictMono.injective h4n

theorem solutions_infinite : solutions.Infinite := by
  apply (Set.infinite_range_of_injective natMap_injective).mono
  rintro _ <n, rfl>
  exact natMap_mem n
\end{lstlisting}

\subsection{Proof status}

\begin{center}
\renewcommand{\arraystretch}{1.3}
\begin{tabular}{lll}
\toprule
Statement & Method & Status \\
\midrule
\texttt{family1 n} & \texttt{ring} & proved \\
\texttt{family2 k} & \texttt{ring} & proved \\
\texttt{family3 t} & \texttt{ring} & proved \\
Explicit witnesses (9 triples) & \texttt{norm\_num} & proved \\
\texttt{natMap\_injective} & \texttt{neg\_inj} + monotonicity & proved \\
\texttt{solutions\_infinite} & injection from $\N$ & proved \\
\midrule
\textbf{Total sorry count} & & \textbf{0} \\
\bottomrule
\end{tabular}
\end{center}

% ════════════════════════════════════════════════════════════════════════════
\section{Computational Verification}\label{sec:computation}

The Python script \texttt{brute\_force\_search.py} performs an exhaustive search over
$|x|, |y| \leq 200$.  For each pair $(x,y)$ it checks whether
$\frac{-x^5 - y^4}{3}$ is a non-negative perfect square, giving $z$.

Within this range the search finds $77$ solutions (counting $\pm z$), all belonging
to one of the three families above or to additional seed solutions that generate
further infinite families.  The 21 primitive seeds found within the range include
$(-1, \pm 1, 0)$, $(-3, 0, \pm 9)$, $(-4, \pm 4, \pm 16)$ as the three simplest.

% ════════════════════════════════════════════════════════════════════════════
\section{Summary}\label{sec:summary}

\begin{center}
\renewcommand{\arraystretch}{1.4}
\begin{tabular}{ll}
\toprule
\multicolumn{2}{c}{$x^5 + y^4 + 3z^2 = 0$} \\
\midrule
Trivial solution & $(0,0,0)$ \\
Seed 1 & $(-1, \pm 1, 0)$ \\
Seed 2 & $(-3, 0, \pm 9)$ \\
Seed 3 & $(-4, \pm 4, \pm 16)$ \\
Family 1 (all $z=0$ solutions) & $(-n^4, \pm n^5, 0)$, $n \in \Z$ \\
Family 2 (all $y=0$ solutions) & $(-3k^2, 0, \pm 9k^5)$, $k \in \Z$ \\
Family 3 & $(-4t^4, \pm 4t^5, \pm 16t^{10})$, $t \in \Z$ \\
Additional primitive seeds ($|x| \leq 200$) & 18 more \\
\midrule
Number of integer solutions & \textbf{Infinite} \\
Lean~4 sorry count & \textbf{0} \\
\bottomrule
\end{tabular}
\end{center}

\end{document}
